\documentclass[a4paper,12pt,titlepage,finall]{article}

\usepackage[T1,T2A]{fontenc}
\usepackage[utf8x]{inputenc}
\usepackage[russian]{babel}
\usepackage{amsmath}
\usepackage{tikz}
\usepackage{pgfplots}
\usepackage{geometry}
\usepackage{indentfirst}
\usepackage{listings}
\usepackage{xcolor}

\geometry{a4paper,left=30mm,top=30mm,bottom=30mm,right=30mm}

\setcounter{secnumdepth}{0}

\usepgfplotslibrary{fillbetween}

\lstset{
    basicstyle=\ttfamily\small,
    breaklines=true,
    frame=single,
    language=C,
    inputencoding=utf8x,
    extendedchars=true,
    keepspaces=true,
}

\begin{document}

% Титульный лист
\begin{titlepage}
    \begin{center}
        {\small \sc Московский государственный университет \\
        имени М.~В.~Ломоносова \\
        Факультет вычислительной математики и кибернетики\\}
        \vfill
        {\Large \sc Отчёт по заданию №6}\\
        ~\\
        {\large \bf <<Сборка многомодульных программ. \\
        Вычисление корней уравнений и определённых интегралов.>>}\\
        ~\\
        {\large \bf Вариант 9 / 1 / 1}
    \end{center}
    \begin{flushright}
        \vfill {Выполнил:\\
        студент 104 группы\\
        Буклов~М.~А.\\
        ~\\
        Преподаватель:\\
        Смирнов~М.~Б.}
    \end{flushright}
    \begin{center}
        \vfill
        {\small Москва\\2026}
    \end{center}
\end{titlepage}

\tableofcontents
\newpage

% ------------------------------------------------------------------
\section{Постановка задачи}

Требуется написать программу, вычисляющую площадь плоской фигуры, ограниченной
тремя кривыми (вариант~9):
\[
    f_1(x) = \frac{3}{(x-1)^2 + 1}, \quad
    f_2(x) = \sqrt{x + 0{,}5}, \quad
    f_3(x) = e^{-x}.
\]

Программа должна:
\begin{enumerate}
    \item найти абсциссы трёх точек попарного пересечения кривых с точностью~$\varepsilon_1$.
          Согласно варианту реализованы \emph{два} метода приближённого решения
          уравнения $f(x) = g(x)$: метод деления отрезка пополам и метод хорд.
          Выбор метода выполняется на этапе компиляции через макрос препроцессора
          \texttt{-DMETHOD\_CHORDS} (по умолчанию --- деление пополам);
    \item вычислить площадь фигуры как алгебраическую сумму определённых
          интегралов методом прямоугольников (средних точек) с точностью~$\varepsilon_2$.
\end{enumerate}

Отрезки для поиска корней определяются аналитически заранее (с анализом смены
знака разности функций). Точности~$\varepsilon_1$ и~$\varepsilon_2$ подбираются
так, чтобы итоговая площадь была вычислена с точностью $\varepsilon = 10^{-3}$.

Функции $f_1$, $f_2$, $f_3$ реализованы на языке ассемблера NASM32 (архитектура IA-32,
соглашение~cdecl). Остальной код написан на языке~Си. Сборка выполняется через
\texttt{make}.

\newpage

% ------------------------------------------------------------------
\section{Математическое обоснование}

\subsection{Анализ функций}

Функция $f_1(x) = 3/((x-1)^2+1)$ является колоколообразной: её максимум $f_1(1)=3$
достигается при $x=1$, а при $x\to\pm\infty$ функция стремится к нулю. Функция
$f_2(x)=\sqrt{x+0{,}5}$ определена при $x\geq -0{,}5$ и монотонно растёт.
Функция~$f_3(x)=e^{-x}$ монотонно убывает, принимая значение~1 при $x=0$.

\subsection{Нахождение отрезков для поиска корней}

Нужно найти три точки пересечения: $f_1\cap f_3$, $f_2\cap f_3$, $f_1\cap f_2$.
Для каждой пары введём разность $\varphi(x) = f_i(x) - f_j(x)$ и найдём отрезок,
на котором она меняет знак.

\textbf{Пересечение $f_1$ и $f_3$} (точка~A). Вычислим значения $f_1 - f_3$ в нескольких
точках:
\begin{align*}
    x = -0{,}5: \quad & f_1 = 0{,}923,\quad f_3 = 1{,}649 \quad\Rightarrow\quad f_1-f_3 = -0{,}726 < 0,\\
    x = 0:      \quad & f_1 = 1{,}500,\quad f_3 = 1{,}000 \quad\Rightarrow\quad f_1-f_3 = +0{,}500 > 0.
\end{align*}
Знак сменился на $[-0{,}5;\; 0]$, поэтому ищем корень на этом отрезке.

\textbf{Пересечение $f_2$ и $f_3$} (точка~B). Вычислим $f_2 - f_3$:
\begin{align*}
    x = 0: \quad & f_2 = 0{,}707,\quad f_3 = 1{,}000 \quad\Rightarrow\quad f_2-f_3 = -0{,}293 < 0,\\
    x = 1: \quad & f_2 = 1{,}225,\quad f_3 = 0{,}368 \quad\Rightarrow\quad f_2-f_3 = +0{,}857 > 0.
\end{align*}
Ищем корень на $[0;\; 1]$.

\textbf{Пересечение $f_1$ и $f_2$} (точка~C). Вычислим $f_1 - f_2$:
\begin{align*}
    x = 1: \quad & f_1 = 3{,}000,\quad f_2 = 1{,}225 \quad\Rightarrow\quad f_1-f_2 = +1{,}775 > 0,\\
    x = 3: \quad & f_1 = 0{,}600,\quad f_2 = 1{,}871 \quad\Rightarrow\quad f_1-f_2 = -1{,}271 < 0.
\end{align*}
Ищем корень на $[1;\; 3]$.

\subsection{Метод деления отрезка пополам}

Пусть $\varphi(x) = f(x) - g(x)$ непрерывна на $[a,b]$ и $\varphi(a)\varphi(b) < 0$.
На каждом шаге вычисляется середина отрезка $c = (a+b)/2$ и значение $\varphi(c)$;
далее один из концов отрезка заменяется на $c$ так, чтобы знак функции на концах
по-прежнему оставался разным. Итерации продолжаются, пока $b - a > \varepsilon_1$,
после чего возвращается $(a+b)/2$.

\subsection{Метод хорд}

Метод хорд (секущих) приближает корень точкой пересечения хорды, соединяющей
концы отрезка~$(a,\varphi(a))$ и~$(b,\varphi(b))$, с осью~$Ox$:
\[
    x_{n+1} = a_n - \varphi(a_n)\,\frac{b_n - a_n}{\varphi(b_n) - \varphi(a_n)}.
\]
Найденную точку~$x_{n+1}$ берут вместо того конца отрезка, на котором знак
функции совпадает со знаком~$\varphi(x_{n+1})$ (так, чтобы корень оставался
внутри новой пары концов). В качестве критерия остановки используется
$|x_{n+1} - x_n| < \varepsilon_1$.

Метод хорд имеет, как правило, более высокую скорость сходимости, чем метод
деления пополам, поскольку использует не только знаки функции на концах, но и
их величины. На тесте $x^2 - 2 = 0$ при $\varepsilon = 10^{-8}$ метод хорд
потребовал~11 итераций против~27 у деления пополам (см. раздел~6).

\subsection{Обоснование выбора $\varepsilon_1$}

Для метода деления пополам на отрезке $[a,b]$ погрешность после $n$ итераций равна~\cite{math}:
\[
    |x_n - x^*| \leq \frac{b-a}{2^n}.
\]
Чтобы гарантировать~$\varepsilon_1 = 10^{-5}$ на самом широком из трёх отрезков
(длина~$[1,3] = 2$), нужно:
\[
    n \geq \log_2\!\frac{2}{10^{-5}} = \log_2(200000) \approx 17{,}6 \;\Rightarrow\; n \geq 18.
\]

Для метода хорд число итераций при той же~$\varepsilon_1$ меньше за счёт более
высокого порядка сходимости. Такое число итераций в обоих случаях выполняется
менее чем за миллисекунду. Поскольку $\varepsilon_1 = 10^{-5} \ll \varepsilon = 10^{-3}$,
погрешность нахождения корней практически не влияет на точность площади.

\subsection{Формула для площади}

На участке $[x_A, x_B]$ порядок функций следующий: $f_2 < f_3 < f_1$
(верхняя граница~--- $f_1$, нижняя~--- $f_3$). На участке $[x_B, x_C]$:
$f_3 < f_2 < f_1$ (нижняя граница~--- $f_2$). Таким образом:
\[
    S = \int_{x_A}^{x_B} \bigl(f_1(x) - f_3(x)\bigr)\,dx
      + \int_{x_B}^{x_C} \bigl(f_1(x) - f_2(x)\bigr)\,dx.
\]

\subsection{Обоснование выбора $\varepsilon_2$}

Погрешность метода средних прямоугольников оценивается как~\cite{math}:
\[
    |I - I_n| \leq \frac{M_2(b-a)^3}{24n^2},
\]
где $M_2 = \max_{[a,b]} |f''(x)|$. В программе реализован критерий Рунге:
\[
    \frac{|I_{2n} - I_n|}{3} < \varepsilon_2.
\]
При $\varepsilon_2 = 10^{-5}$ фактическая погрешность каждого интеграла
не превышает $10^{-5}$, что в сумме по двум слагаемым даёт не более
$2 \cdot 10^{-5} \ll 10^{-3}$.

\subsection{Графики функций}

\begin{figure}[h]
\centering
\begin{tikzpicture}
\begin{axis}[
    axis lines=middle,
    restrict x to domain=-1.5:3,
    restrict y to domain=-0.5:4,
    axis equal,
    enlargelimits,
    legend cell align=left,
    scale=1.8]

\addplot[green!70!black,samples=300,thick] {3/((x-1)^2+1)};
\addlegendentry{$y=\dfrac{3}{(x-1)^2+1}$}

\addlegendimage{empty legend}\addlegendentry{}

\addplot[blue,domain=-0.5:3,samples=300,thick] {sqrt(x+0.5)};
\addlegendentry{$y=\sqrt{x+0{,}5}$}

\addplot[red,samples=300,thick] {exp(-x)};
\addlegendentry{$y=e^{-x}$}

\end{axis}
\end{tikzpicture}
\caption{Графики функций $f_1$, $f_2$, $f_3$}
\label{plot1}
\end{figure}

\newpage

% ------------------------------------------------------------------
\section{Результаты экспериментов}

Программа запускалась с $\varepsilon_1 = \varepsilon_2 = 10^{-5}$.

\begin{table}[h]
\centering
\begin{tabular}{|c|c|c|}
\hline
Пересечение & $x$ & $y$ \\
\hline
$f_1$ и $f_3$ (точка A) & $-0{,}2033$ & $1{,}2255$ \\
$f_2$ и $f_3$ (точка B) & $\phantom{-}0{,}1874$ & $0{,}8291$ \\
$f_1$ и $f_2$ (точка C) & $\phantom{-}1{,}9562$ & $1{,}5672$ \\
\hline
\end{tabular}
\caption{Координаты вершин фигуры}
\label{table1}
\end{table}

Площадь фигуры:
\[
    S_{AB} = \int_{-0{,}2033}^{0{,}1874} (f_1 - f_3)\,dx \approx 0{,}18877, \quad
    S_{BC} = \int_{0{,}1874}^{1{,}9562} (f_1 - f_2)\,dx \approx 2{,}14982,
\]
\[
    S = S_{AB} + S_{BC} \approx \mathbf{2{,}3386}.
\]

\begin{figure}[h]
\centering
\begin{tikzpicture}
\begin{axis}[
    axis lines=middle,
    restrict x to domain=-1.5:3,
    restrict y to domain=-0.5:4,
    axis equal,
    enlargelimits,
    legend cell align=left,
    scale=1.8,
    xticklabels={,,},
    yticklabels={,,}]

\addplot[green!70!black,samples=300,thick,name path=F1] {3/((x-1)^2+1)};
\addlegendentry{$f_1$}

\addlegendimage{empty legend}\addlegendentry{}

\addplot[blue,domain=-0.5:3,samples=300,thick,name path=F2] {sqrt(x+0.5)};
\addlegendentry{$f_2$}

\addplot[red,samples=300,thick,name path=F3] {exp(-x)};
\addlegendentry{$f_3$}

\addplot[blue!25,samples=300] fill between[of=F1 and F3, soft clip={domain=-0.2033:0.1874}];
\addplot[blue!25,samples=300] fill between[of=F1 and F2, soft clip={domain=0.1874:1.9562}];
\addlegendentry{$S \approx 2{,}3386$}

\addplot[dashed,gray] coordinates {(-0.2033, 1.2255) (-0.2033, 0)};
\addplot[color=black] coordinates {(-0.2033, 0)} node [label={-135:{\small $-0{,}2033$}}]{};

\addplot[dashed,gray] coordinates {(0.1874, 0.8291) (0.1874, 0)};
\addplot[color=black] coordinates {(0.1874, 0)} node [label={-45:{\small $0{,}1874$}}]{};

\addplot[dashed,gray] coordinates {(1.9562, 1.5672) (1.9562, 0)};
\addplot[color=black] coordinates {(1.9562, 0)} node [label={-90:{\small $1{,}9562$}}]{};

\end{axis}
\end{tikzpicture}
\caption{Закрашенная фигура}
\label{plot2}
\end{figure}

\newpage

% ------------------------------------------------------------------
\section{Структура программы и спецификация функций}

Программа состоит из шести модулей.

\begin{itemize}
    \item \textbf{f1.asm}~--- вычисление $f_1(x) = 3/((x-1)^2+1)$ средствами x87 FPU.
    \item \textbf{f2.asm}~--- вычисление $f_2(x) = \sqrt{x+0{,}5}$ инструкцией \texttt{fsqrt}.
    \item \textbf{f3.asm}~--- вычисление $f_3(x) = e^{-x}$ через тождество
          $e^y = 2^{y \log_2 e}$ с помощью \texttt{fldl2e}, \texttt{f2xm1}, \texttt{fscale}.
    \item \textbf{root.c}~--- функция \texttt{root(f, g, a, b, eps)} для уравнения
          $f(x)=g(x)$. Содержит две реализации: метод деления отрезка пополам
          (выбирается по умолчанию) и метод хорд (выбирается через макрос
          \texttt{METHOD\_CHORDS}, переключение --- директивами \texttt{\#ifdef}).
          Хранит число итераций в глобальной переменной \texttt{g\_iterations}.
    \item \textbf{integral.c}~--- функция \texttt{integral(f, a, b, eps)}: метод средних
          прямоугольников с остановкой по критерию Рунге.
    \item \textbf{main.c}~--- разбор аргументов командной строки, вычисление площади,
          режимы тестирования.
\end{itemize}

Все функции $f_1$, $f_2$, $f_3$ и их вызывающие стороны следуют соглашению \textbf{cdecl}:
аргумент \texttt{double} передаётся на стеке по адресу \texttt{[ebp+8]}, результат
возвращается в регистре \texttt{st0} сопроцессора x87.

Зависимости между модулями показаны на рис.~\ref{dep}.

\begin{figure}[h]
\centering
\begin{tikzpicture}[
    box/.style={draw, minimum width=2.2cm, minimum height=0.8cm, align=center},
    arr/.style={->, thick}]

\node[box] (main)     at (0,0)    {main.c};
\node[box] (root)     at (-3,-2)  {root.c};
\node[box] (integral) at ( 3,-2)  {integral.c};
\node[box] (f1)       at (-5,-4)  {f1.asm};
\node[box] (f2)       at (-3,-4)  {f2.asm};
\node[box] (f3)       at (-1,-4)  {f3.asm};
\node[box] (f1r)      at ( 1,-4)  {f1.asm};
\node[box] (f2r)      at ( 3,-4)  {f2.asm};
\node[box] (f3r)      at ( 5,-4)  {f3.asm};
\node[box] (funcs)    at (0,-6)   {funcs.h};

\draw[arr] (main) -- (root);
\draw[arr] (main) -- (integral);
\draw[arr] (root) -- (f1);
\draw[arr] (root) -- (f2);
\draw[arr] (root) -- (f3);
\draw[arr] (integral) -- (f1r);
\draw[arr] (integral) -- (f2r);
\draw[arr] (integral) -- (f3r);

\end{tikzpicture}
\caption{Зависимости между модулями}
\label{dep}
\end{figure}

\newpage

% ------------------------------------------------------------------
\section{Сборка программы (Make-файл)}

Сборка производится командой \texttt{make}. Ниже приведён Makefile:

\begin{lstlisting}[language=make,basicstyle=\ttfamily\small]
CC       = gcc
NASM     = nasm
CFLAGS   = -m32 -std=c99
NFLAGS   = -f elf32
OBJECTS  = main.o root.o integral.o f1.o f2.o f3.o

ifeq ($(METHOD),CHORDS)
    CFLAGS += -DMETHOD_CHORDS
endif

all: prog

prog: $(OBJECTS)
	$(CC) $(CFLAGS) -o prog $(OBJECTS) -lm

main.o:     src/main.c include/funcs.h
	$(CC) $(CFLAGS) -c -o main.o src/main.c

root.o:     src/root.c include/funcs.h
	$(CC) $(CFLAGS) -c -o root.o src/root.c

integral.o: src/integral.c include/funcs.h
	$(CC) $(CFLAGS) -c -o integral.o src/integral.c

f1.o: asm/f1.asm
	$(NASM) $(NFLAGS) -o f1.o asm/f1.asm

f2.o: asm/f2.asm
	$(NASM) $(NFLAGS) -o f2.o asm/f2.asm

f3.o: asm/f3.asm
	$(NASM) $(NFLAGS) -o f3.o asm/f3.asm

clean:
	rm -f *.o prog
\end{lstlisting}

Объектные файлы \texttt{f1.o}, \texttt{f2.o}, \texttt{f3.o} собираются NASM в формате
\texttt{elf32}. Объектные файлы Си-модулей собираются GCC с флагом~\texttt{-m32} для
32-разрядного кода. Финальная сборка выполняется gcc с подключением \texttt{libm} для
функций стандартной математической библиотеки, используемых в~\texttt{integral.c}.
Выбор метода поиска корня осуществляется переменной \texttt{METHOD}: по умолчанию
собирается версия с делением пополам, команда \texttt{make METHOD=CHORDS} добавляет
к \texttt{CFLAGS} флаг \texttt{-DMETHOD\_CHORDS}, и компилируется ветка с методом хорд.

\newpage

% ------------------------------------------------------------------
\section{Отладка программы, тестирование функций}

\subsection{Тестирование \texttt{root}}

Функция \texttt{root} тестируется командой \texttt{./prog --test-root}.

\textbf{Тест~1.} Нахождение $\sqrt{2}$ как корня $x^2 - 2 = 0$ на $[1, 2]$.
Аналитический ответ: $x^* = \sqrt{2} \approx 1{,}41421356$. Функция $\varphi(x) = x^2 - 2$:
$\varphi(1) = -1 < 0$, $\varphi(2) = 2 > 0$ --- знак меняется, оба метода применимы.
Для деления пополам теоретическое число итераций:
\[
    n \geq \log_2\frac{2-1}{10^{-8}} \approx 27 \text{ итераций.}
\]
Результаты программы при $\varepsilon = 10^{-8}$:
\begin{itemize}
    \item деление пополам (\texttt{make}): $x \approx 1{,}4142135642$, итераций~27;
    \item метод хорд (\texttt{make METHOD=CHORDS}): $x \approx 1{,}4142135605$, итераций~11.
\end{itemize}
Оба значения имеют ошибку $< 10^{-8}$. Метод хорд сошёлся примерно в~2{,}5 раза
быстрее за счёт более высокого порядка сходимости.

\textbf{Тест~2.} Нахождение корня $\cos(x) = 0$ на $[1, 2]$ через \texttt{root(cos\_f, zero, 1, 2, 1e-7)}.
Аналитический ответ: $x^* = \pi/2 \approx 1{,}5707963$. Знак: $\cos(1) \approx 0{,}54 > 0$,
$\cos(2) \approx -0{,}41 < 0$ --- применим.

\textbf{Тест~3.} Нахождение корня $e^x - 2 = 0$ на $[0, 1]$.
Аналитический ответ: $x^* = \ln 2 \approx 0{,}693147$. Знак:
$e^0 - 2 = -1 < 0$, $e^1 - 2 \approx 0{,}718 > 0$ --- применим.

\subsection{Тестирование \texttt{integral}}

Функция \texttt{integral} тестируется командой \texttt{./prog --test-integral}.

\textbf{Тест~1.} $\displaystyle\int_0^1 x\,dx = 0{,}5$ (прямая, тривиальный). Результат: $0{,}5000000000$.

\textbf{Тест~2.} $\displaystyle\int_0^1 x^2\,dx = \tfrac{1}{3} \approx 0{,}333333$.
Функция $x^2$ --- квадратичная, метод средних прямоугольников интегрирует её точно
за одно разбиение (нет погрешности за счёт симметрии). Результат: $0{,}3333333333$.

\textbf{Тест~3.} $\displaystyle\int_0^\pi \sin(x)\,dx = 2$. Функция нелинейная, интеграл
вычисляется с несколькими разбиениями. Результат: $2{,}0000000000$ (до 10 знаков).

\newpage

% ------------------------------------------------------------------
\section{Программа на Си и на Ассемблере}

Исходные тексты программы приложены к отчёту в архиве. Архив содержит:
\begin{itemize}
    \item \texttt{src/main.c}, \texttt{src/root.c}, \texttt{src/integral.c} --- код на Си;
    \item \texttt{asm/f1.asm}, \texttt{asm/f2.asm}, \texttt{asm/f3.asm} --- ассемблерные
          реализации функций варианта;
    \item \texttt{include/funcs.h} --- заголовочный файл с объявлениями всех функций;
    \item \texttt{Makefile} --- сценарий сборки.
\end{itemize}

\newpage

% ------------------------------------------------------------------
\section{Анализ допущенных ошибок}

В ходе разработки были допущены следующие ошибки.

\textbf{1. Неверные границы поиска корней.} Изначально отрезок для $f_1 \cap f_2$ был взят
как $[-2, 0]$, не учитывая, что $f_2(x) = \sqrt{x+0{,}5}$ не определена при $x < -0{,}5$.
Это приводило к вычислению $\sqrt{\text{отрицательного числа}}$ и возврату \texttt{NaN}.
Ошибка устранена аналитической проверкой области определения функций до выбора отрезков.

\textbf{2. Стек x87 после \texttt{fscale}.} В реализации $f_3$ инструкция \texttt{fscale}
оставляет на стеке два значения: результат и исходное~$n$. Без явного удаления~$n$
(инструкция \texttt{fstp st1}) стек FPU переполнялся после нескольких вызовов функции,
и результаты становились некорректными. После добавления \texttt{fstp st1} ошибка исчезла.

\textbf{3. Кеш Fa в методе деления пополам.} В первоначальной реализации значение
$F(a) = f(a) - g(a)$ вычислялось заново на каждой итерации, хотя правый конец отрезка
обновляется только при $F_a \cdot F_{mid} < 0$, а иначе меняется только $a$. После
кеширования~$F_a$ число вызовов функций сократилось примерно вдвое.

\newpage

\begin{raggedright}
\addcontentsline{toc}{section}{Список цитируемой литературы}
\begin{thebibliography}{99}
\bibitem{math} Ильин~В.\,А., Садовничий~В.\,А., Сендов~Бл.\,Х.
    Математический анализ. Т.\,1~--- Москва: Наука, 1985.
\end{thebibliography}
\end{raggedright}

\end{document}
